\documentclass[12pt]{article}

% Packages
\usepackage[margin=1in]{geometry}
\usepackage{amsmath}
\usepackage{amssymb}
\usepackage{amsthm}
\usepackage{enumitem}

% Title information
\title{Mathematics 629\\
Spring 2026\\
Homework Assignment 6}
\author{Alexander Kim}
\date{Due: 3/25/2026}

\begin{document}

\maketitle

\begin{enumerate}

	\item Write out the complete details of the proof of the ``differentiation under the integral'' Theorem~1.129 in the notes.

	      \begin{proof}
		      i) Suppose there exists $g \in L^1(\mu)$ such that $|f(x,t)| \leq g(x)$ for all $(x,t)$ and $\lim_{t\to t_o} f(x,t) = f(x,t_o)$. Let $f_n(x) = f(x, t_o)$ so $f_n(x) \to f(x,t)$. Since $g(x)$ dominates $f(x,t)$ and $f_n$ converges pointwise, by DCT,
		      $$\lim \int f_n d\mu = \int f(x,t_o) d\mu$$
		      or equivalently $F(t_n)\to F(t_o)$. Since $t_n \to t_o$ is an arbitrary sequence then
		      $\lim_{t\to t_o} F(t) = F(t_o)$ for all $t_o$.
		      \newline
		      ii) Suppose there also exists $t\in [a,b]$ and partial derivative $\frac{\partial f}{\partial t}$ exists such that $|\frac{\partial f}{\partial t} f(x,t)| \leq g(x)$ almost everywhere.
		      Let $d_n(x) = \frac{f(x,t + h_n - f(x, t_o)}{h_n}$ be a sequence where $h_n \to 0$. Then by definition of derivative, $d_n\to \frac{\partial f}{\partial t}$. By assumption, $|d_n| \leq |\frac{\partial f}{\partial t}| \leq g(x)$ a.e. so by DCT,
		      $$\lim_{n\to \infty} \int d_n d\mu = \int \frac{\partial f}{\partial t}(x,t) d\mu.$$
		      Since $\int d_n d\mu = \int  \frac{f(x,t + h_n) - f(x, t_o)}{h_n} = \frac{F(t) - F(t_o)}{h_n}$ then $F'(t) \to \int \frac{\partial f}{\partial t}(x,t) d\mu.$
	      \end{proof}

	\item Prove that $\int_{(0,\infty)} x^n e^{-x}\,dx = n!$ by differentiating $\int_0^\infty e^{-tx}\,dt = 1/x$.

	      \begin{proof}
		      Starting with $\int_0^\infty e^{-tx} dt = \frac{1}{x}$, differentiate both sides such that
		      $\frac{d^n}{dx^n} \frac{1}{x} = \frac{(-1)^n n!}{x^{m+1}}$ and
		      $\frac{d^n}{dx^n} \int_0^\infy e^{-tx}dt = \int_0^\infty \frac{d^n}{dx^n}(e^{-tx})dt = \int_0^\infty(-1)^n t^n e^{-tx}dt$. Then $\int_0^\infty t^n e^{-tx}dt = \frac{n!}{x^{n+1}}$ so when $x = 1$, $\int_0^\infty t^ne^{-tx}dt = n!$. If $x \in [a,b]$ then $t^ne^{-t} \leq t^ne^{-tx} \leq t^n e^{-ta}$so the necessary conditions of thoerem 1.129 are satisfied and the differentiation under the integral sign is valid.
	      \end{proof}

	\item Prove that if $f_n$ is a sequence of decreasing nonnegative measurable functions with $\int f_1\,d\mu < \infty$, then
	      \[
		      \int \lim_{n\to\infty} f_n\,d\mu = \lim_{n\to\infty} \int f_n\,d\mu.
	      \]

	      \begin{proof}
		      Since $\int f_1 d\mu < \infty$ then $f_1 \in L^1(\mu)$ and $f_1$ dominates $f_n$ as a sequence of decreasing nonnegative measurable functions. Furthermore, $\lim_{n\to \infty} f_n = 0$ since $f_n$ is bounded below by 0 and decreasing. Thus DCT shows
		      $\int \lim_{n\to \infty} f_n d\mu = \lim_{n\to \infty}\int f_n d\mu.$
	      \end{proof}

	\item Let $y_1, \ldots, y_M$ be distinct points in $\mathbb{R}$, and $\beta_1, \ldots, \beta_M \in \mathbb{R}$. Let
	      \[
		      f(x) = \prod_{j=1}^{M} |x - y_j|^{-\beta_j}.
	      \]
	      Find a necessary and sufficient condition on $(\beta_1, \ldots, \beta_M)$ for $f$ to belong to $L^1(\mathbb{R}, m)$.

	      \begin{proof}
		      The necessary and sufficient condition for $f\in L^1(\mathbb{R},m)$ is that $\beta_j < 1$ for all $j\in \{1, 2, \cdots M\}$ and $\sum_{j=1}^M \beta_j < 1$ also. $\beta_j < 1$ is necessary to prevent $f(x)$ from diverging near $y_j$ but this isn't sufficient because a small $\beta_j$ like $\beta_j = \frac{1/2M}$ would cause $\int f dm$ to diverge. Enforcing $\sum_{j=1}^M \beta_j <1$ ensures $\int f dm$ is finite since $f(x) = |x - y_j|^{\beta_{product} < 1}$ whose integral does converge.
	      \end{proof}

	\item Define $g\colon \mathbb{R} \to \mathbb{R}$ by
	      \[
		      g(x) = \begin{cases} |x|^{-1/3}, & |x| < 1, \\ 0, & |x| > 1. \end{cases}
	      \]
	      Let $\{r_k\}_{k=0}^\infty$ be an enumeration of all points in $[0,1] \cap \mathbb{Q}$. Let
	      \[
		      f(x) = \mathbf{1}_{[0,1]} \sum_{k \in \mathbb{N}} 2^{-k}\, g(x - r_k).
	      \]

	      \begin{enumerate}[label=(\roman*)]
		      \item Show that $f \in L^1([0,1])$.

		            \begin{proof}
			            Let $f_N(x) = \mathbf{1}_{[0,1]}\sum_{k=1}^N 2^{-k} g(x - r_k)dx$ be the Nth partial sum. It follows that
			            \begin{align}
				            \int_0^1 f(x)dx & = \sum_{k=1}^N 2^{-k} \int_0^1 g(x)dx                                         \\
				                            & = \sum_{k=1}^N \int_{-r_k}^0 (r_k-x)^{-1/3}dx + \int_0^{r_k} (x-r_k)^{-1/3}dx \\
				                            & = \frac{3}{2}r_k^{2/3} + \frac{3}{2}(1-r_k)^{2/3} \leq 3.
			            \end{align}
			            Since $f_N(x) \to f(x)$ and $\int_0^1 f(x)dx$ is bounded then by MCT $\int_0^1 f(x)dx = \lim_{N\to \infty}\int_0^1 f_N(x)dx \leq 3$ so $f$ is $\mathcal{L}^1([0,1])$ and thus $L^1([0,1])$ also.
		            \end{proof}

		      \item Show that $f$ is unbounded on every open set (and remains so after modification on a set of measure zero). Show that $f$ is discontinuous at every point in $[0,1]$ and remains so after modification on a set of measure zero.

		            \begin{proof}
			            Let $O\subset [0,1]$ be any subset of $[0,1]$. Since $\mathbb{Q}$ is dense then $r_k \in O$ and when
			            $x\neq r_k$, $g(x-r_k) = |x-r_k|^{-1/3}$. Notice that $f(x) \geq \frac{|x-r_k|^{-1/3}}{2^k} \to \infty$ as $r_k \to x$ so $f$ is unbounded on $O$. Now let $h=f$ a.e. such that
			            $E = \{x: h(x) \neq f(x)\}$ has $m(E) = 0$. For any $C>0$
			            $$\{x\in O: f(x) >M\} \subseteq \{x\in O:f(x)>M\}\backslash M = \{x\in O: h(x) >M\}$$
			            which has non-zero measure, so $h(x)$ is also unbounded on $O$. Thus, $f$ is unbounded on open intervals regardless of modification by a nullset.

		            \end{proof}

		      \item Let $0 < \varepsilon < 1$. Construct a compact subset $K_\varepsilon$ of $[0,1]$ such that $m(K_\varepsilon) \geq 1 - \varepsilon$ and $f\mathbf{1}_{K_\varepsilon}$ is continuous on $K_\varepsilon$.

		            \begin{proof}
			            Let $0<\epsilon<1$ and let $O_k = (r_k- \frac{\epsilon}{2^{k+2}}, r_k + \frac{\epsilon}{2^{k+2}})$.
			            Observe that given $U = \cup_{k\in \mathbb{N}}O_k$, $m(U)\leq \sum_{k\in \mathbb{N}}m(O_k) = \frac{\epsilon}{2}$ by subadditivity.
			            Define $K_\epsilon = [0,1]\backslash U$ and observe that $m(K_\epsilon)\geq 1-\frac{\epsilon}{2} \geq 1 - \epsilon.$ Then for all $x \in K_\epsilon$, $|x-r_k| \geq \frac{\epsilon}{2^{k+1}}$ since
			            $x\notin U$. Let $f_N(x) = \sum_{k=1}^N 2^{-k}g(x-r_k)$ and see
			            $$sup|f(x) - f_N(x)| \leq \sum_{k=N+1}^\infty 2^{-k} (\frac{\epsilon}{2^{k+2}})^{-1/3} = \frac{2^{2/3}}{\epsilon^{1/3}}\sum_{k=N+1}^\infty 2^{-2k/3}\to 0$$
			            by ratio test of geometric series. Thus $f$ is uniformly continuous on $K_\epsilon$.

		            \end{proof}
	      \end{enumerate}

	\item Let $f \in L^1([0, 2\pi])$.

	      \begin{enumerate}[label=(\roman*)]
		      \item Prove that $\displaystyle\lim_{n\to\infty} \int_0^{2\pi} f(x)\,e^{-inx}\,dx = 0$.

		            \begin{proof}
			            First observe that $|e^{-inx}| \leq 1$. By theorem 1.117 let $g$ be a step function such that $\int_0^{2\pi} |f - g| < \epsilon$. Then by the triangle inequality,
			            $$|\int_0^{2\pi} f e^{-inx}dx| \leq |\int_0^{2\pi} (f-g)e^{-inx}dx| + |\int_0^{2\pi} ge^{-inx}dx|.$$
			            Since $g$ is a step function given sufficiently large $n$ we know $|\int_0^{2\pi}ge^{-inx}|\to 0$. For the other piece, notice
			            $$|\int_0^{2\pi} (f-g)e^{-inx}dx| \leq \int |f-g|dx < \epsilon.$$
			            Thus by substituting, $\int_0^{2\pi}f(x)^e^{-inx}dx \leq \lim|\int_0^{2\pi}f(x)e^{-inx}dx| < \epsilon$ for sufficiently large $n$ so $\lim_{n\to\infty}fe^{-inx}dx = 0$.
		            \end{proof}

		      \item Prove that for a set $E \subset [0, 2\pi]$ of positive Lebesgue measure $m(E)$,
		            \[
			            \lim_{n\to\infty} \int_E \cos^2(nx + \gamma_n)\,dx = \frac{m(E)}{2},
		            \]
		            for any sequence $\gamma_n$.

		            \begin{proof}
			            To begin use the trig identity for $cos^2(\theta)$ such that $cos^2(nx+\gamma_n) = \frac{1 + cos(2nx+2\gamma_n)}{2}.$ Then $\int_E cos^2(nx+\gamma_n)dx = \frac{m(E)}{2} + \frac{1}{2} \int_E cos(2nx + 2\gamma_n)dx.$ Using euler's formula notice $cos(2nx+2\gamma_n) = Re(e^{2nx}e^{2\gamm_n})$ so by substitution $\int_E cos(2nx + 2\gamma_n)dx = Re(e^{2\gamma_n}\int_E e^{2nx}dx) \leq |\int e^{2inx}dx|$ since $|e^{2i\gamma_n}| = 1$. By part i),
			            $$\lim_{n\to\infy}\int_E e^{2inx}dx = 0$$ so
			            $$\int_E cos(2nx+2\gamma_n)dx \to 0$$
			            since $E\subset [0,2\pi]$.
			            By substitution, $\lim_{n\to\infy}\int_E cos^2(nx+\gamma_n)dx = \frac{m(E)}{2}.$
		            \end{proof}

		      \item Let $F_n(x) = a_n \cos nx + b_n \sin nx$. Assume that the series $\sum_{n=0}^\infty F_n(x)$ converges on a set $E$ of positive measure. Show that $\lim_{n\to\infty} a_n = 0$ and $\lim_{n\to\infty} b_n = 0$.

		            \textit{Hint: It might help to express $F_n(x)^2$ in the form $(r_n \cos(nx + \gamma_n))^2$, with $r_n > 0$.}

		            \begin{proof}
			            Let $F_n(x)^2 = r^2cos(nx+\gamma_n)^2$ where $r^2 = a_n^2+b_n^2$. We want to show $r_n\to 0$. Since $\sum_{n=0}^\infty F_n(x)$ converges on $E$ then for all $x\in E$, $F_n(x)\to 0$ and furthermore $F_n(x)^2\to 0$ pointwise also. Using Egorov's theorem, there exists a subset $A\subset E$ such that $F_n^2(x) \to 0$ uniformly on $A$ which shows $\int_A F_n^2(x)dx \to 0$.
			            Since $\int_A F_n(x)^2dx = r_n^2\int_A cos(nx+\gamma_n)^2dx$ then
			            $$r_n^2 = \frac{\int_A F_n(x)^2dx}{\int_A cos(nx+\gamma_n)^2dx}.$$
			            Using the result part ii)
			            $$r_n^2 \to \frac{0}{m(A)/2} = 0$$
			            thus $a_n^2+b_n^2 \to 0$ so individually $a_n\to 0$ and $b_n \to 0$.
		            \end{proof}
	      \end{enumerate}

\end{enumerate}

\end{document}
